% v2.4.0 P3 figure batch 02: supervised, EM, covariance, variational, and PageRank figures.
\documentclass[UTF8,a4paper,11pt,openany]{ctexbook}
\usepackage{../common/statlearnbook}
\SLSetBookMetadata{P3 绘图批次 02}{统计学习方法讲义 v2.4.0 绘图集成验证}{第二册至第五册局部编译}
\begin{document}
\pagestyle{empty}
\raggedbottom
\typeout{SLFIG-HARNESS-BEGIN FIG-P222-01}
% v2.3.1 figure UID: FIG-P222-01
% v2.6.0 Image-reference reverse redraw: clarify the class-to-feature path without changing topology.
\tikzset{
  slfig-FIG-P222-01/.style={font=\fontsize{9.2pt}{11.0pt}\selectfont,every path/.append style={line cap=round,line join=round}},
  slfig-FIG-P222-01-var/.style={circle,draw=SLBlue,line width=.82pt,minimum size=7.5mm,inner sep=0pt,fill=white},
  slfig-FIG-P222-01-class/.style={slfig-FIG-P222-01-var,fill=SLBlueBG,line width=.98pt},
  slfig-FIG-P222-01-arrow/.style={-{Stealth[length=2.0mm,width=1.25mm]},draw=SLBlue,line width=.92pt},
  slfig-FIG-P222-01-category/.style={font=\fontsize{9.2pt}{11.0pt}\selectfont,text=SLBlue},
  slfig-FIG-P222-01-region/.style={draw=SLTeal,fill=SLTealBG,line width=.66pt},
  slfig-FIG-P222-01-region-label/.style={font=\fontsize{9.2pt}{11.0pt}\selectfont,text=SLTeal},
  slfig-FIG-P222-01-formula/.style={slfig note,font=\fontsize{9.2pt}{11.0pt}\selectfont,
    align=center,rounded corners=1.4mm,inner xsep=2.5mm,inner ysep=1.6mm,line width=.60pt}
}
\begin{figure}[H]
\centering
\begin{tikzpicture}[slfig-FIG-P222-01]
  \node[slfig-FIG-P222-01-class] (y) at (0,2.25) {$Y$};
  \node[slfig-FIG-P222-01-category,above=2mm of y] {类别变量};
  \node[slfig-FIG-P222-01-var] (x1) at (-3.50,0) {$X_1$};
  \node[slfig-FIG-P222-01-var] (x2) at (-1.75,0) {$X_2$};
  \node at (0,0) {$\cdots$};
  \node[slfig-FIG-P222-01-var] (xm) at (1.75,0) {$X_{d-1}$};
  \node[slfig-FIG-P222-01-var] (xd) at (3.50,0) {$X_d$};
  \draw[slfig-FIG-P222-01-arrow] (y)--(x1);
  \draw[slfig-FIG-P222-01-arrow] (y)--(x2);
  \draw[slfig-FIG-P222-01-arrow] (y)--(xm);
  \draw[slfig-FIG-P222-01-arrow] (y)--(xd);
  \begin{scope}[on background layer]
  \node[slfig-FIG-P222-01-region,rounded corners=2mm,fit=(x1)(x2)(xm)(xd),
        inner xsep=5mm,inner ysep=7mm] {};
  \end{scope}
  \node[slfig-FIG-P222-01-region-label] at (0,-.58) {给定 $Y$ 后相互独立};
  \node[slfig-FIG-P222-01-formula,anchor=north] at (0,-1.28)
    {$X_i\perp X_j\mid Y\quad(i\ne j)$};
\end{tikzpicture}
\caption{朴素贝叶斯的条件依赖结构：给定$Y$后，各特征节点相互独立}\label{fig:V2-C03-star}
\end{figure}
\clearpage
\typeout{SLFIG-HARNESS-END FIG-P222-01}
\typeout{SLFIG-HARNESS-BEGIN FIG-P262-01}
% v2.6.0 figure UID: FIG-P262-01
% Image-reference reverse redraw: preserve the semantic geometry and clarify the visual hierarchy.
\tikzset{
  slfig-FIG-P262-01/.style={
    every path/.append style={line cap=round,line join=round},
    font=\fontsize{9.2}{11}\selectfont,
  },
  slfig-FIG-P262-01-direct/.style={
    slfig direct label,font=\fontsize{9.2}{11}\selectfont,inner sep=1.1pt,
  },
  slfig-FIG-P262-01-note/.style={
    slfig note,font=\fontsize{9.2}{11}\selectfont,inner xsep=3.2pt,inner ysep=2.2pt,
  },
}
\pgfplotsset{
  slfig-FIG-P262-01-axis/.style={
    axis line style={draw=SLInk!82,line width=.52pt},
    tick style={draw=SLInk!76,line width=.44pt},
    tick label style={font=\fontsize{8.7}{10.2}\selectfont},
    label style={font=\fontsize{9.5}{11.2}\selectfont},
    clip=false,
  },
  slfig-FIG-P262-01-curve/.style={draw=SLBlue,line width=1.02pt},
  slfig-FIG-P262-01-reference/.style={draw=SLInk!38,line width=.52pt,dash pattern=on 4.2pt off 2.5pt},
  slfig-FIG-P262-01-tangent/.style={draw=SLGold!82!black,line width=.78pt},
  slfig-FIG-P262-01-guide/.style={draw=SLInk!48,line width=.52pt,dash pattern=on 2.2pt off 1.8pt},
}
\begin{figure}[H]
\centering
\begin{tikzpicture}[slfig-FIG-P262-01]
\begin{axis}[slfig axis,slfig-FIG-P262-01-axis,
  width=.82\linewidth,height=.50\linewidth,
  xmin=-6.25,xmax=6.25,ymin=-.08,ymax=1.08,
  axis lines=middle,xlabel={$z$},ylabel={$\sigma(z)$},
  xtick={-2,0,2},xticklabels={$-a$,0,$a$},
  ytick={0,.5,1},yticklabels={0,$\tfrac12$,1},
  clip=false,
]
  \addplot[slfig-FIG-P262-01-curve,domain=-6:6,samples=301]
    {1/(1+exp(-x))};
  \addplot[slfig-FIG-P262-01-reference] coordinates {(-6,0) (6,0)};
  \addplot[slfig-FIG-P262-01-reference] coordinates {(-6,1) (6,1)};
  \addplot[slfig-FIG-P262-01-tangent,domain=-1.55:1.55,samples=2]
    {.5+.25*x};
  \addplot[slfig-FIG-P262-01-guide] coordinates {(-2,0) (-2,{1/(1+exp(2))})};
  \addplot[slfig-FIG-P262-01-guide] coordinates {(2,0) (2,{1/(1+exp(-2))})};
  \addplot[only marks,mark=*,mark size=2.2pt,draw=SLTeal,fill=SLTeal]
    coordinates {(-2,{1/(1+exp(2))}) (2,{1/(1+exp(-2))})};
  \addplot[only marks,mark=*,mark size=2.3pt,draw=SLGold!82!black,fill=SLGold]
    coordinates {(0,.5)};
  \node[slfig-FIG-P262-01-direct,text=SLBlue,anchor=west] at (axis cs:2.65,.82) {概率映射};
  \draw[draw=SLGold!82!black,line width=.46pt]
    (axis cs:.55,.6375)--(axis cs:.82,.50);
  \node[font=\fontsize{9.2}{11}\selectfont,text=SLGold!82!black,anchor=west]
    at (axis cs:.95,.40) {中心斜率 $\sigma'(0)=\tfrac14$};
  \node[slfig-FIG-P262-01-note,anchor=east] at (axis cs:-2.55,.72)
    {$\sigma(-a)=1-\sigma(a)$};
\end{axis}
\end{tikzpicture}
\caption{逻辑斯谛函数把线性预测量映射为概率，并满足关于$(0,1/2)$的中心对称性}\label{fig:V2-C05-sigmoid}
\end{figure}
\clearpage
\typeout{SLFIG-HARNESS-END FIG-P262-01}
\typeout{SLFIG-HARNESS-BEGIN FIG-P346-01}
% v2.3.1 figure UID: FIG-P346-01
% Proposed post-v2.3.1 polish for FIG-P346-01: protect key-label size and stroke hierarchy.
\tikzset{slfig-FIG-P346-01/.style={every path/.append style={line cap=round,line join=round}}}
\pgfplotsset{slfig-FIG-P346-01-axis/.style={
  tick label style={font=\fontsize{8.5pt}{10pt}\selectfont},
  label style={font=\fontsize{9.5pt}{11pt}\selectfont},clip=false}}
\begin{figure}[H]
\centering
\begin{tikzpicture}[slfig-FIG-P346-01]
\begin{axis}[slfig-FIG-P346-01-axis,
  name=boundaxis,
  slfig axis,width=.82\linewidth,height=.51\linewidth,
  xmin=.2,xmax=6.2,ymin=-.3,ymax=4.35,
  axis lines=left,xlabel={$\theta$},ylabel={目标值},
  xtick={2},xticklabels={2},ytick=\empty,clip=false,
]
  \addplot[name path=ell,draw=SLBlue,line width=1.05pt,domain=.3:6.1,samples=321]
    {3.8-.18*(x-4.5)^2};
  \addplot[name path=bound,draw=SLTeal,line width=.88pt,densely dashed,domain=.3:6.1,samples=321]
    {3.8-.18*(x-4.5)^2-.16*(x-2)^2};
  \addplot[draw=none,fill=SLTeal,fill opacity=.055]
    fill between[of=ell and bound];
  \addplot[slfig reference,line width=.52pt] coordinates {(2,-.05) (2,2.675)};
  \addplot[draw=SLGold,line width=.76pt,domain=1.22:2.90,samples=2]
    {2.675+.9*(x-2)};
  \addplot[only marks,mark=*,mark size=2.5pt,draw=SLGold,fill=SLGold]
    coordinates {(2,2.675)};
  \draw[draw=SLBlue,line width=.46pt]
    (axis cs:4.70,3.793)--(axis cs:4.94,4.04);
  \node[slfig direct label,font=\fontsize{9.2pt}{10.8pt}\selectfont,text=SLBlue,
    fill=none,inner sep=0pt,anchor=west] at (axis cs:5.04,4.10) {$\ell(\theta)$};
  \draw[draw=SLTeal,line width=.46pt]
    (axis cs:4.50,2.80)--(axis cs:4.57,2.47);
  \node[slfig direct label,font=\fontsize{9.2pt}{10.8pt}\selectfont,text=SLTeal,
    fill=none,inner sep=0pt,anchor=west] at (axis cs:4.65,2.34) {$B(\theta,2)$：下界};
  \draw[draw=SLGold!88!black,line width=.46pt]
    (axis cs:2.00,2.675)--(axis cs:1.58,3.16);
  \node[font=\fontsize{9pt}{10.5pt}\selectfont,text=SLGold!88!black,
    anchor=west] at (axis cs:.62,3.38) {相切点；共享切线};
\end{axis}
  \node[slfig note,font=\fontsize{9pt}{10.8pt}\selectfont,line width=.58pt,
    fill=white,anchor=north,inner xsep=4mm,inner ysep=2.2mm]
    at ([yshift=-10mm]boundaxis.south)
    {$\ell(\theta)-B(\theta,2)=0.16(\theta-2)^2\ge0$};
\end{tikzpicture}
\captionsetup{width=.94\linewidth}
\caption{可复算的相切下界：$B(\theta,2)=\ell(\theta)-0.16(\theta-2)^2\le\ell(\theta)$，并在$\theta=2$处函数值与导数同时相等。两条曲线按图中给定的显式函数精确绘制，而非未标比例的定性示意。}\label{fig:V3-C04-bound}
\end{figure}
\clearpage
\typeout{SLFIG-HARNESS-END FIG-P346-01}
\typeout{SLFIG-HARNESS-BEGIN FIG-P482-01}
% v2.3.1 figure UID: FIG-P482-01
% Proposed post-v2.3.1 polish for FIG-P482-01: protect key-label size and stroke hierarchy.
\tikzset{
  slfig-FIG-P482-01/.style={font=\fontsize{9.4pt}{11.2pt}\selectfont,every path/.append style={line cap=round,line join=round}},
  slfig-FIG-P482-01-residual/.style={draw=SLRed,line width=.76pt,densely dashed},
  slfig-FIG-P482-01-right-angle/.style={draw=SLRed,line width=.62pt},
}
\pgfplotsset{
  slfig-FIG-P482-01-axis/.style={
    axis line style={draw=SLGray!88,line width=.58pt},
    tick label style={font=\fontsize{8.8pt}{10.5pt}\selectfont},
    label style={font=\fontsize{10pt}{12pt}\selectfont},clip=false,
  },
  slfig-FIG-P482-01-samples/.style={only marks,mark=*,mark size=1.55pt,
    draw=SLBlue!72,fill=SLBlue!18,line width=.48pt,fill opacity=.72},
  slfig-FIG-P482-01-inner/.style={draw=SLTeal!72,line width=.58pt},
  slfig-FIG-P482-01-outer/.style={draw=SLTeal,line width=.82pt},
  slfig-FIG-P482-01-axis-one/.style={draw=SLBlue,line width=1.08pt},
  slfig-FIG-P482-01-axis-two/.style={draw=SLGray,line width=.72pt},
  slfig-FIG-P482-01-mean/.style={only marks,mark=*,mark size=2.55pt,draw=SLGold,fill=SLGold},
  slfig-FIG-P482-01-query/.style={only marks,mark=triangle*,mark size=2.9pt,
    draw=SLRed,fill=white,line width=.85pt},
  slfig-FIG-P482-01-projection/.style={only marks,mark=square*,mark size=2pt,
    draw=SLRed,fill=SLRed},
}
\begin{figure}[H]
\centering
\begin{tikzpicture}[slfig-FIG-P482-01]
\begin{axis}[slfig-FIG-P482-01-axis,
  slfig axis,width=.78\linewidth,height=.61\linewidth,
  xmin=-3.4,xmax=3.8,ymin=-2.45,ymax=2.55,
  axis equal image,axis lines=middle,xlabel={$x_1$},ylabel={$x_2$},
  xtick=\empty,ytick=\empty,clip=false,
]
  \addplot[slfig-FIG-P482-01-samples]
    coordinates {(-2.45,-1.18) (-2.05,-.42) (-1.72,-1.02) (-1.45,-.05)
      (-1.05,-.62) (-.72,.30) (-.35,-.35) (.02,.58) (.38,.12) (.72,.82)
      (1.05,.32) (1.34,1.08) (1.66,.52) (1.98,1.32) (2.28,.78) (2.72,1.46)
      (-.92,-1.12) (.38,-.48) (1.18,-.08) (2.10,.18)};
  % Covariance eigenpairs: lambda_1=2.25, lambda_2=.36, angle=25 deg.
  \addplot[slfig-FIG-P482-01-inner,domain=0:360,samples=241]
    ({.3+1.3595*cos(x)-.2536*sin(x)},{-.1+.6339*cos(x)+.5438*sin(x)});
  \addplot[slfig-FIG-P482-01-outer,domain=0:360,samples=241]
    ({.3+2.7189*cos(x)-.5071*sin(x)},{-.1+1.2679*cos(x)+1.0876*sin(x)});
  \addplot[slfig-FIG-P482-01-axis-one]
    coordinates {(-2.4189,-1.3679) (3.0189,1.1679)};
  \addplot[slfig-FIG-P482-01-axis-two]
    coordinates {(.8071,-1.1876) (-.2071,.9876)};
  \addplot[slfig-FIG-P482-01-mean]
    coordinates {(.3,-.1)};
  \node[font=\fontsize{9.4pt}{11.2pt}\selectfont,text=SLGold,anchor=north east]
    at (axis cs:-.08,-.52) {均值 $\mu$};
  \node[slfig direct label,text=SLBlue,anchor=west] at (axis cs:2.78,1.33) {第一主轴 $\lambda_1$};
  \node[slfig direct label,text=SLGray,anchor=east] at (axis cs:-.30,1.22) {第二主轴 $\lambda_2$};
  \node[font=\fontsize{9.4pt}{11.2pt}\selectfont,text=SLTeal,anchor=north east]
    at (axis cs:-2.15,-1.52) {$2\sigma$};
  \node[font=\fontsize{9.4pt}{11.2pt}\selectfont,text=SLTeal!78,anchor=north]
    at (axis cs:-.96,-.96) {$1\sigma$};

  % q = mu + 1.2 e_1 + .7 e_2; its projection is mu + 1.2 e_1.
  \addplot[slfig-FIG-P482-01-query]
    coordinates {(1.0918,1.0415)};
  \draw[slfig-FIG-P482-01-residual]
    (axis cs:1.0918,1.0415)--(axis cs:1.3876,.4071);
  \addplot[slfig-FIG-P482-01-projection]
    coordinates {(1.3876,.4071)};
  \draw[slfig-FIG-P482-01-right-angle]
    (axis cs:1.337,.385)--(axis cs:1.359,.334)--(axis cs:1.410,.356);
  \node[slfig direct label,text=SLRed,anchor=south west] at (axis cs:1.16,1.12) {样本 $q$};
  \node[slfig direct label,text=SLRed,anchor=west] at (axis cs:1.48,.31) {正交投影};
\end{axis}
\end{tikzpicture}
\caption{二维协方差椭圆与主轴示意}\label{fig:V4-C04-ellipse}
\end{figure}
\clearpage
\typeout{SLFIG-HARNESS-END FIG-P482-01}
\typeout{SLFIG-HARNESS-BEGIN FIG-P689-01}
% v2.6.0 figure UID: FIG-P689-01
% Iteration 05: use an explicit 8 mm x-shift to move the complete right-hand axis into its panel.
\tikzset{slfig-FIG-P689-01/.style={font=\fontsize{9.2pt}{11.0pt}\selectfont,every path/.append style={line cap=round,line join=round}}}
\pgfplotsset{slfig-FIG-P689-01-axis/.style={tick label style={font=\footnotesize},label style={font=\small},clip=false}}
\pgfkeysifdefined{/tikz/alt}{}{\tikzset{alt/.code={}}}
\begin{figure}[htbp]
\centering
\begin{tikzpicture}[slfig-FIG-P689-01,
  every node/.style={font=\fontsize{9.2pt}{11.0pt}\selectfont,align=center,outer sep=0pt},
  panel/.style={draw=SLRuleGray,line width=.55pt,rounded corners=3pt,
    minimum width=67mm,minimum height=53mm,inner sep=0pt},
  alt={对象—关系—结论：左侧以长度分解直接显示观测对数证据等于ELBO与非负KL间隙之和；右侧坐标更新使ELBO形成非降阶梯，但只能到坐标稳定或局部驻点，未知全局上限不应被解释为已经达到。}]

\node[panel] at (-3.85,0.10) {};
\node[font=\bfseries,text=SLDeepBlue] at (-3.85,2.42)
  {证据的长度分解};
\fill[SLMidBlue,opacity=0.08] (-6.55,0.50) rectangle (-2.55,1.25);
\fill[fill=SLAccentOrange,opacity=0.08] (-2.55,0.50) rectangle (-0.95,1.25);
\draw[draw=SLDeepBlue,line width=.95pt] (-6.55,0.50) rectangle (-0.95,1.25);
\draw[draw=SLMidBlue,line width=.90pt] (-2.55,0.50) -- (-2.55,1.25);
\node[text=SLDeepBlue] at (-4.55,0.875) {$\mathcal L(q)$：证据下界};
\node[text=SLAccentOrange] at (-1.75,0.875) {KL\ 间隙};
\draw[<->,draw=SLTextGray,line width=.58pt] (-6.55,1.46) -- (-0.95,1.46)
  node[midway,above=2.0mm] {$\log p(w)$};
\node[align=left,text width=59mm] at (-3.85,-0.15)
  {$\displaystyle \log p(w)=\mathcal L(q)+
    \operatorname{KL}\!\left(q(h)\,\|\,p(h\mid w)\right)$};
\node[align=left,text width=59mm,text=SLTextGray] at (-3.85,-1.35)
  {KL$\ge0$，故$\mathcal L(q)\le\log p(w)$；变分族受限时，间隙可保持为正。};

\node[panel] at (3.85,0.10) {};
\node[font=\bfseries,text=SLDeepBlue] at (3.85,2.42)
  {坐标更新下的 ELBO 非降阶梯};
\begin{axis}[slfig-FIG-P689-01-axis,
  at={(0.80,-1.55)},anchor=south west,xshift=8mm,yshift=-10.5mm,
  width=59mm,height=34mm,scale only axis,
  xmin=0,xmax=6.3,ymin=0,ymax=1.12,
  axis lines=left,axis line style={draw=SLTextGray,line width=.55pt},
  tick style={draw=SLRuleGray},xtick={0,1,2,3,4,5,6},ytick=\empty,
  xlabel={坐标更新轮次},xlabel style={font=\fontsize{9.0pt}{10.6pt}\selectfont},
  clip=false]
  \addplot+[const plot,draw=SLMidBlue,line width=.95pt,mark=*,mark size=1.5pt]
    coordinates {(0,0.12) (1,0.34) (2,0.50) (3,0.64) (4,0.73) (5,0.79) (6,0.80)};
  \addplot+[draw=SLAccentOrange,dashed,line width=.70pt,mark=none]
    coordinates {(0,1.02) (6.3,1.02)};
  \node[anchor=north west,font=\fontsize{9.0pt}{10.6pt}\selectfont,text=SLAccentOrange]
    at (axis cs:0.15,0.98) {未知全局上限};
  \node[anchor=north east,font=\fontsize{9.0pt}{10.6pt}\selectfont,text=SLDeepBlue]
    at (axis cs:6.18,0.48) {坐标稳定／局部驻点};
\end{axis}
\end{tikzpicture}
\captionsetup{width=.94\linewidth}
\caption{观测对数证据等于ELBO与变分分布到真实后验的KL散度之和；平均场坐标更新可使ELBO逐步不降，但非凸目标下有限运行通常只得到坐标稳定点或局部驻点，多启动比较也不构成全局最优证明}
\label{fig:V5-C06-elbo-geometry}
\end{figure}
\clearpage
\typeout{SLFIG-HARNESS-END FIG-P689-01}
\typeout{SLFIG-HARNESS-BEGIN FIG-P690-01}
% v2.6.0 figure UID: FIG-P690-01
% Mean-field comparison: real formula bands, an unobstructed transition corridor,
% and a geometric break instead of an opaque white eraser.
\tikzset{slfig-FIG-P690-01/.style={font=\fontsize{9.2pt}{11.0pt}\selectfont,every path/.append style={line cap=round,line join=round}}}
\pgfkeysifdefined{/tikz/alt}{}{\tikzset{alt/.code={}}}
\begin{figure}[htbp]
\centering
\begin{tikzpicture}[slfig-FIG-P690-01,
  >=Stealth,
  every node/.style={font=\fontsize{9.2pt}{11.0pt}\selectfont,align=center},
  latent/.style={circle,draw=SLMidBlue,fill=white,minimum size=10mm,inner sep=1pt},
  qlatent/.style={latent,minimum size=15.5mm,font=\fontsize{9.2pt}{11.0pt}\selectfont},
  observed/.style={circle,draw=SLDeepBlue,fill=SLDeepBlue,text=white,minimum size=10mm,inner sep=1pt},
  param/.style={draw=SLRuleGray,fill=SLSoftGray,rounded corners=2pt,minimum width=20mm,minimum height=8mm},
  panel/.style={draw=SLRuleGray,rounded corners=3pt,densely dashed,inner xsep=5.8mm,inner ysep=5.5mm},
  dep/.style={draw=SLMidBlue,line width=0.90pt},
  input/.style={->,draw=SLDeepBlue,line width=0.80pt},
  cutedge/.style={draw=SLRuleGray,densely dashed,line width=0.55pt},
  cutmark/.style={draw=SLAccentOrange,line width=0.70pt},
  >={Stealth[length=2.2mm,width=1.6mm]},
  alt={对象—关系—结论：左右面板使用完全相同的节点位置。真实单文档后验中主题比例与主题指派耦合；平均场近似切断这一后验依赖，但仍保留观测词与固定主题词参数对责任度更新的输入。}]

\node[font=\bfseries\fontsize{9.6pt}{11.5pt}\selectfont,text=SLDeepBlue] at (-4.40,2.55)
  {真实后验：$\theta_m$ 与 $z_{mn}$ 耦合};
\node[latent] (tl) at (-6.40,0.75) {$\theta_m$};
\node[latent] (zl) at (-4.40,0.75) {$z_{mn}$};
\node[observed] (wl) at (-2.40,0.75) {$w_{mn}$};
\node[param] (pl) at (-4.40,-0.85) {$\boldsymbol\varphi$ 固定};
\draw[dep] (tl) -- (zl);
\draw[input] (wl) -- (zl);
\draw[input] (pl) -- (zl);
\node[panel,fit=(tl)(zl)(wl)(pl)] (leftpanel) {};
\node[text=SLTextGray] at (-4.40,-2.25)
  {$p(\theta_m,z_{mn}\mid w_m,\boldsymbol\varphi)$};

\node[font=\bfseries\fontsize{9.6pt}{11.5pt}\selectfont,text=SLDeepBlue] at (4.40,2.55)
  {平均场：切断近似后验依赖};
\node[qlatent] (tr) at (2.40,0.75) {$q(\theta_m)$};
\node[qlatent] (zr) at (4.40,0.75) {$q(z_{mn})$};
\node[observed] (wr) at (6.40,0.75) {$w_{mn}$};
\node[param] (pr) at (4.40,-0.85) {$\boldsymbol\varphi$ 固定};
\draw[cutedge] (tr.east) -- (3.21,0.75);
\draw[cutedge] (3.59,0.75) -- (zr.west);
\draw[cutmark] (3.34,0.59) -- (3.39,0.89);
\draw[cutmark] (3.41,0.59) -- (3.46,0.89);
\draw[input] (wr) -- (zr);
\draw[input] (pr) -- (zr);
\node[panel,fit=(tr)(zr)(wr)(pr)] (rightpanel) {};
\node[text=SLTextGray] at (4.40,-2.25)
  {$q(\theta_m)\prod_n q(z_{mn})$};

\draw[->,draw=SLDeepBlue,line width=0.85pt] (leftpanel.east) -- (rightpanel.west)
  node[midway,above=2.2mm,font=\fontsize{9.2pt}{10.8pt}\selectfont,
       text width=26mm,align=center]{选择可计算的\\乘积分布族};
\node[draw=SLMidBlue,fill=SLSoftBlue,rounded corners=2pt,
  text width=125mm,minimum height=10mm,inner xsep=3.0mm,inner ysep=1.8mm] at (0,-3.84)
  {切断的是近似后验中的直接依赖，不是删除生成模型依赖；\\
   $w_{mn}$ 与$\boldsymbol\varphi$仍进入$q(z_{mn})$的责任度更新。};
\end{tikzpicture}
\caption{固定主题--词参数$\boldsymbol\varphi$时，平均场近似把单文档后验中耦合的$\boldsymbol\theta_m$与$z_{mn}$替换为两个变分因子族；被切断的是近似后验中的直接依赖，固定参数仍通过词似然进入责任度更新}
\label{fig:V5-C06-mean-field-graph}
\end{figure}
\clearpage
\typeout{SLFIG-HARNESS-END FIG-P690-01}
\typeout{SLFIG-HARNESS-BEGIN FIG-P717-01}
% v2.7.0 figure UID: FIG-P717-01
% Basic/no-dangling inbound balance, with the general-PageRank boundary local.
\tikzset{slfig-FIG-P717-01/.style={
  font=\fontsize{9.6pt}{11.5pt}\selectfont,
  every path/.append style={line cap=round,line join=round}}}
\pgfkeysifdefined{/tikz/alt}{}{\tikzset{alt/.code={}}}
\begin{figure}[htbp]
\centering
\begin{tikzpicture}[slfig-FIG-P717-01,>=Stealth,
  every node/.style={
    font=\fontsize{9.6pt}{11.5pt}\selectfont,align=center},
  source/.style={
    circle,draw=SLRuleGray,line width=.74pt,fill=SLSoftGray,
    minimum size=11mm,inner sep=0pt},
  target/.style={
    circle,draw=SLDeepBlue,line width=.96pt,fill=SLDeepBlue,
    text=white,minimum size=15mm,inner sep=0pt},
  inbound/.style={
    -{Stealth[length=2.4mm,width=1.8mm]},
    draw=SLMidBlue,line width=.96pt},
  aggregate/.style={
    -{Stealth[length=2.4mm,width=1.8mm]},
    draw=SLDeepBlue,line width=1.0pt},
  contribution/.style={
    font=\fontsize{9.6pt}{11.5pt}\selectfont,
    text=SLDeepBlue,fill=white,rounded corners=1pt,
    inner sep=1.0pt},
  formula/.style={
    draw=SLDeepBlue,line width=.82pt,fill=white,
    rounded corners=3pt,text width=62mm,minimum height=35mm,
    inner xsep=2.5mm,inner ysep=2.3mm,align=left},
  boundary/.style={
    draw=SLTextGray,dashed,line width=.82pt,fill=SLWarmGray,
    rounded corners=3pt,text width=138mm,
    inner xsep=2.5mm,inner ysep=2.0mm,align=left},
  alt={对象—关系—结论：上部严格限定为基本PageRank且无悬挂结点，此处修复符号 S 等于链接矩阵 M。三个来源 j1、j2、j3 以完全相同线宽指向目标 i，每条边只表示第 j 列提供贡献 M 下标 ij 乘 r_j 的上一步值，不编码数值大小；目标再把贡献汇总。公式卡说明非加权时存在 j 到 i 当且仅当 M_ij 等于来源 j 的出度倒数，分母属于来源 j。下部边界卡说明一般PageRank先以 S 等于 M 加 v a 转置修复悬挂列；悬挂列即使没有 j 到 i 也可有 S_ij 等于 v_i 大于零，随后 G 等于 dS 加一减 d 乘 v 乘全一向量转置，所以正元素对应入链不能外推到修复后的 S 或 G。}]

\node[
  font=\fontsize{10.2pt}{12.2pt}\selectfont\bfseries,
  text=SLDeepBlue] at (0,3.10)
  {基本 PageRank 入链贡献（无悬挂；此处 $S=M$）};

% Three source columns feed one target with exactly equal edge widths.
\node[source] (j1) at (-6.25,1.30) {$j_1$};
\node[source] (j2) at (-6.25,0) {$j_2$};
\node[source] (j3) at (-6.25,-1.30) {$j_3$};
\node[target] (i) at (-2.15,0) {目标 $i$};

\draw[inbound] (j1) to[out=0,in=150]
  node[contribution,pos=.42,above=2.4mm]
  {$M_{i j_1}r_{j_1}^{(t)}$} (i);
\draw[inbound] (j2) --
  node[contribution,pos=.42,above=2.0mm]
  {$M_{i j_2}r_{j_2}^{(t)}$} (i);
\draw[inbound] (j3) to[out=0,in=210]
  node[contribution,pos=.42,below=2.4mm]
  {$M_{i j_3}r_{j_3}^{(t)}$} (i);

\node[formula] (formula) at (3.20,0)
  {\textbf{来源 $j$ 的第 $j$ 列 $\longrightarrow$ 目标 $i$}\\
   $\displaystyle r_i^{(t+1)}=\sum_j M_{ij}r_j^{(t)}$\\
   $\displaystyle \hphantom{r_i^{(t+1)}}
   =\sum_{j:j\to i}M_{ij}r_j^{(t)}$\\
   基本 PageRank、无悬挂：$S=M$，$M_{ij}=A_{ij}/c_j$。\\
   非加权时：$M_{ij}=1/\deg^+(j)\iff j\to i$；\\
   否则 $M_{ij}=0$。分母 $\deg^+(j)$ 属于\textbf{来源 $j$}。\\
   三条入边等线宽；未给数值，不编码贡献大小。};

\draw[aggregate] (i) -- (formula)
  node[midway,above=1.8mm,fill=white,inner sep=1pt,
  text=SLDeepBlue] {汇总};

% The boundary travels with the figure, so the following section title cannot
% visually extend the basic/no-dangling claim to repaired S or to G.
\node[boundary] at (0,-4.20)
  {\textbf{一般 PageRank 边界（不可外推上方“正元素对应入链”）}\\
   先以 $S=M+\boldsymbol v\boldsymbol a^{\mathsf T}$ 修复悬挂列；若
   $c_j=0$，即使没有 $j\to i$，也可有 $S_{ij}=v_i>0$。\\
   随后 $G=dS+(1-d)\boldsymbol v\boldsymbol1^{\mathsf T}$；因此上方入链解释
   只属于基本无悬挂的 $M(=S)$，不能套到修复后的 $S/G$。};

\end{tikzpicture}
\captionsetup{width=.92\linewidth}
\caption{基本 PageRank（无悬挂）的入链贡献及一般 PageRank 的适用边界。}
\label{fig:V5-C07-inbound-contribution}
\end{figure}
\clearpage
\typeout{SLFIG-HARNESS-END FIG-P717-01}
\typeout{SLFIG-HARNESS-BEGIN FIG-P720-01}
% v2.3.1 figure UID: FIG-P720-01
% v2.6.0 reverse redraw for FIG-P720-01: preserve all 19 coordinates and 15 directed iterates while opening the stationary-point and formula-card layout.
\tikzset{slfig-FIG-P720-01/.style={font=\fontsize{9.2pt}{11.0pt}\selectfont,every path/.append style={line cap=round,line join=round}}}
\pgfkeysifdefined{/tikz/alt}{}{\tikzset{alt/.code={}}}
\begin{figure}[htbp]
\centering
\begin{tikzpicture}[slfig-FIG-P720-01,
  >=Stealth,
  every node/.style={font=\fontsize{9.2pt}{11.0pt}\selectfont,align=center},
  patha/.style={draw=SLDeepBlue,line width=0.80pt},
  pathb/.style={draw=SLMidBlue,line width=0.80pt,dashed},
  pathc/.style={draw=SLTextGray,line width=0.78pt,dash dot},
  alt={对象—关系—结论：三个真实初值在同一三状态概率单纯形上按固定PageRank算子迭代；三条路径的步长逐次缩短并收敛到唯一星形平稳点，同心收缩区与一范数不等式共同说明压缩率不超过d。}]
\coordinate (e1) at (-5.8,-1.70); \coordinate (e2) at (-1.2,-1.70); \coordinate (e3) at (-3.5,2.28);
\fill[fill=SLSoftBlue,opacity=.04] (e1)--(e2)--(e3)--cycle;
\draw[draw=SLDeepBlue,line width=0.94pt] (e1)--(e2)--(e3)--cycle;
\node[below=1.2mm] at (e1) {$e_1$}; \node[below=1.2mm] at (e2) {$e_2$}; \node[above=1.2mm] at (e3) {$e_3$};
\coordinate (star) at (-3.313,0.339);
\draw[draw=SLRuleGray,line width=0.50pt,fill=SLSoftBlue,fill opacity=.030] (star) circle (0.78);
\draw[draw=SLRuleGray,line width=0.50pt,fill=SLSoftBlue,fill opacity=.030] (star) circle (0.42);

\coordinate (a0) at (-5.110,-1.302); \coordinate (a1) at (-2.764,0.263); \coordinate (a2) at (-3.414,0.554); \coordinate (a3) at (-3.318,0.232); \coordinate (a4) at (-3.300,0.373);
\coordinate (b0) at (-1.890,-1.302); \coordinate (b1) at (-3.408,1.378); \coordinate (b2) at (-3.414,-0.041); \coordinate (b3) at (-3.249,0.430); \coordinate (b4) at (-3.336,0.330);
\coordinate (c0) at (-3.500,1.484); \coordinate (c1) at (-3.408,-0.108); \coordinate (c2) at (-3.242,0.455); \coordinate (c3) at (-3.341,0.324); \coordinate (c4) at (-3.306,0.334);
\foreach \s/\t/\al/\aw in {a0/a1/1.65mm/1.18mm,a1/a2/1.40mm/1.00mm,a2/a3/1.05mm/.78mm,a3/a4/.65mm/.48mm,a4/star/.35mm/.26mm}{\draw[patha,-{Stealth[length=\al,width=\aw]}] (\s)--(\t);}
\foreach \s/\t/\al/\aw in {b0/b1/1.65mm/1.18mm,b1/b2/1.45mm/1.04mm,b2/b3/1.10mm/.80mm,b3/b4/.60mm/.44mm,b4/star/.30mm/.23mm}{\draw[pathb,-{Stealth[length=\al,width=\aw]}] (\s)--(\t);}
\foreach \s/\t/\al/\aw in {c0/c1/1.55mm/1.10mm,c1/c2/1.25mm/.90mm,c2/c3/.72mm/.52mm,c3/c4/.35mm/.26mm,c4/star/.20mm/.16mm}{\draw[pathc,-{Stealth[length=\al,width=\aw]}] (\s)--(\t);}
\node[draw=SLDeepBlue,rectangle,minimum size=3.4mm,inner sep=0pt] at (a0) {};
\node[draw=SLMidBlue,circle,minimum size=3.4mm,inner sep=0pt] at (b0) {};
\node[draw=SLTextGray,regular polygon,regular polygon sides=3,minimum size=4.2mm,inner sep=0pt] at (c0) {};
\node[star,star points=5,draw=SLDeepBlue,fill=SLDeepBlue,text=white,minimum size=6.0mm,inner sep=0pt] at (star) {};
\node[text=white,inner sep=0pt] at (star) {$r^*$};
\node[text=SLTextGray,font=\fontsize{8.8pt}{10.5pt}\selectfont] at (-3.5,-2.56)
  {三条路径均由同一$G$的真实数值迭代生成};

\node[draw=SLRuleGray,rounded corners=3pt,fill=white,text width=78mm,minimum height=42mm,
  inner xsep=2.0mm,inner ysep=2.0mm,align=left] at (4.00,0.16)
  {$T(x)=Gx=dSx+(1-d)v$，$0\le d<1$。\\[4.5pt]
   $\displaystyle \|Gx-Gy\|_1\le d\|x-y\|_1$\\[4.5pt]
   $\displaystyle \|r^{(t)}-r^*\|_1\le d^t\|r^{(0)}-r^*\|_1$。\\[4.5pt]
   初值形状不同，终点唯一；同心区域表示到$r^*$的距离持续收缩。};
\end{tikzpicture}
\caption{阻尼PageRank在概率单纯形上的压缩迭代与唯一平稳点}
\label{fig:V5-C07-simplex-contraction}
\end{figure}
\clearpage
\typeout{SLFIG-HARNESS-END FIG-P720-01}
\typeout{SLFIG-HARNESS-BEGIN FIG-P721-01}
% v2.7.0 figure UID: FIG-P721-01
% v2.7.0: exact six-decimal trajectory, non-stopping t=6 display cut,
%          and separately computed exact fixed point/ranking.
\tikzset{slfig-FIG-P721-01/.style={font=\fontsize{9.6pt}{11.5pt}\selectfont,every path/.append style={line cap=round,line join=round}}}
\pgfkeysifdefined{/tikz/alt}{}{\tikzset{alt/.code={}}}
\begin{figure}[H]
\centering
\begin{tikzpicture}[slfig-FIG-P721-01,
  every node/.style={font=\fontsize{9.6pt}{11.5pt}\selectfont},
  every pin/.style={font=\fontsize{9.6pt}{11.5pt}\selectfont},
  alt={对象—关系—结论：左面板给出从均匀初值独立复算的 t=0 到 8 三个排名分量有限迭代轨迹；t=6 的橙色竖线只是展示截断，图内非零固定点残差明确排除停止判定。右面板与有限轨迹分开，给出精确固定点 r 星的六位小数和排名 3 大于 2 大于 1；三条轨迹同时用颜色、线型和点型编码。}]
\begin{scope}[shift={(-6.8,-2.25)}]
\begin{groupplot}[
  group style={group size=2 by 1,horizontal sep=20mm},
  width=58mm,height=47mm,scale only axis,
  axis line style={draw=SLTextGray,line width=.65pt},
  tick style={draw=SLRuleGray,line width=.52pt},
  title style={font=\fontsize{10.2pt}{12.2pt}\selectfont},
  label style={font=\fontsize{9.6pt}{11.5pt}\selectfont},
  tick label style={font=\fontsize{9.6pt}{11.5pt}\selectfont}]
\nextgroupplot[title={迭代轨迹},xlabel={$t$},ylabel={$r_i^{(t)}$},
  xmin=0,xmax=8,ymin=.10,ymax=.56,xtick={0,...,8},ytick={.2,.3,.4,.5},
  ymajorgrids=true,xmajorgrids=false,grid style={draw=SLRuleGray,dotted,line width=.46pt},clip=false]
\addplot[draw=SLDeepBlue,line width=.94pt,mark=*,mark size=1.55pt,mark options={solid,fill=SLDeepBlue}] coordinates {(0,.333333)(1,.155556)(2,.214815)(3,.202173)(4,.202594)(5,.203718)(6,.203074)(7,.203297)(8,.203247)};
\addplot[draw=SLMidBlue,line width=.92pt,dashed,mark=square*,mark size=1.50pt,mark options={solid,fill=SLMidBlue}] coordinates {(0,.333333)(1,.288889)(2,.277037)(3,.288099)(4,.283463)(5,.284756)(6,.284561)(7,.284527)(8,.284566)};
\addplot[draw=SLTextGray,line width=.90pt,dash dot,mark=triangle*,mark size=1.75pt,mark options={solid,fill=SLTextGray}] coordinates {(0,.333333)(1,.555556)(2,.508148)(3,.509728)(4,.513942)(5,.511526)(6,.512365)(7,.512176)(8,.512187)};
\draw[draw=SLAccentOrange,dashed,line width=.74pt] (axis cs:6,.10)--(axis cs:6,.56);
\node[anchor=west,text=SLAccentOrange,align=left,fill=white,
  fill opacity=.94,text opacity=1,rounded corners=1pt,inner sep=1.1pt]
  at (axis cs:6.12,.410) {展示截断 $t=6$\\（非停止判定）};
\node[anchor=south east,text=SLAccentOrange,align=right,fill=white,
  fill opacity=.94,text opacity=1,rounded corners=1pt,inner sep=1.1pt]
  at (axis cs:5.82,.108)
  {$\lVert G\boldsymbol r^{(6)}-\boldsymbol r^{(6)}\rVert_1$\\
   $=4.474947\!\times\!10^{-4}>0$};
\node[inner sep=0pt,pin={[pin distance=5.2mm,pin edge={draw=SLDeepBlue,line width=.42pt},text=SLDeepBlue]150:{网页1}}] at (axis cs:5,.2037) {};
\node[inner sep=0pt,pin={[pin distance=5.2mm,pin edge={draw=SLMidBlue,line width=.42pt},text=SLMidBlue]120:{网页2}}] at (axis cs:5,.2848) {};
\node[inner sep=0pt,pin={[pin distance=5.2mm,pin edge={draw=SLTextGray,line width=.42pt},text=SLTextGray]240:{网页3}}] at (axis cs:5,.5115) {};
\nextgroupplot[title={精确固定点与排名},xlabel={$r_i^*$},xbar,xmin=0,xmax=.58,
  symbolic y coords={网页1,网页2,网页3},ytick={网页1,网页2,网页3},enlarge y limits=.23,
  bar width=4mm,bar shift=0pt,grid=none,clip=false]
\addplot[fill=SLDeepBlue,fill opacity=.75,draw=SLDeepBlue,postaction={pattern=north east lines,pattern color=white}] coordinates {(.203252,网页1)};
\addplot[fill=SLMidBlue,fill opacity=.75,draw=SLMidBlue,postaction={pattern=dots,pattern color=white}] coordinates {(.284553,网页2)};
\addplot[fill=SLTextGray,fill opacity=.75,draw=SLTextGray] coordinates {(.512195,网页3)};
\node[anchor=west,inner sep=0pt] at (axis cs:.235,网页1) {0.203252（第3）};
\node[anchor=west,inner sep=0pt] at (axis cs:.318,网页2) {0.284553（第2）};
\node[anchor=center,text=SLTextGray,inner sep=0pt,
  font=\fontsize{9.6pt}{11.5pt}\selectfont\bfseries,yshift=-5.2mm]
  at (axis cs:.40,网页3) {0.512195（第1）};
\end{groupplot}
\node at (6.8,-1.20) {实线/圆：网页1\quad 虚线/方：网页2\quad 点划线/三角：网页3};
\end{scope}
\end{tikzpicture}
\caption{PageRank 的有限迭代轨迹与精确固定点（显示六位小数）：
$t=6$ 仅为展示截断且残差非零；固定点排序为 $3\succ2\succ1$。}
\label{fig:V5-C07-rank-trajectory}
\end{figure}
\clearpage
\typeout{SLFIG-HARNESS-END FIG-P721-01}
\end{document}
